\section{Modelo de Contenido (UWE) y Modelo Relacional de Base de Datos}
\label{sec:diagrama_clases_detallado}
\label{sec:diseno_bbdd}
% =============================================================================

En la metodología UWE (Koch et al., 2008), el \textbf{Modelo de Contenido} representa mediante un diagrama de clases UML los objetos de información que administra la aplicación web, identificados con el estereotipo «contentClass». En el sistema electoral, este modelo define tanto las clases en memoria con las dimensiones exactas de las primitivas post-cuánticas como el esquema relacional en Tercera Forma Normal (3FN) que desvincula la identidad del ciudadano respecto a su boleta cifrada.

\subsection{Diagrama de Clases de Contenido UWE y Actores Criptográficos}

En la Figura \ref{fig:uml_clases_detallado} se presenta el Diagrama de Clases UML del Modelo de Contenido UWE junto con las cuatro entidades participantes del protocolo, especificando atributos, métodos y multiplicidades según el estándar UML.

\begin{figure}[!htbp]
	\centering
	\resizebox{\textwidth}{!}{%
	\begin{tikzpicture}[
		cls/.style={rectangle, draw=black, fill=white, line width=0.8pt, rounded corners=1.5pt, font=\tiny\rmfamily, align=left, inner sep=3.5pt},
		assoc/.style={-{Latex[length=2mm, width=1.4mm]}, line width=0.8pt, draw=black},
		comp/.style={{Diamond[length=2.5mm, fill=black]}-, line width=0.8pt, draw=black}
	]
		\node[cls] (V) at (-6.2, 4.0) {
			\textbf{VotantePQC}\\
			\rule{4.6cm}{0.3pt}\\
			- id\_v: str\\
			- sk\_v: bytes [4032]\\
			- pk\_v: bytes [1952]\\
			- credencial: CredencialVotante\\
			\rule{4.6cm}{0.3pt}\\
			+ formular\_voto(candidato: str, id\_elec: str): PaqueteVoto\\
			+ encapsular\_clave(ek\_ac: bytes): Tuple[bytes, bytes]\\
			+ firmar\_boleta(boleta: BoletaSufragio): bytes
		};

		\node[cls] (AR) at (-6.2, -0.6) {
			\textbf{AutoridadRegistroPQC}\\
			\rule{4.6cm}{0.3pt}\\
			- nombre: str\\
			- sk\_ar: bytes [4032]\\
			- pk\_ar: bytes [1952]\\
			\rule{4.6cm}{0.3pt}\\
			+ registrar\_ciudadano(id\_v: str, pk\_v: bytes): CredencialVotante\\
			+ verificar\_identidad(id\_v: str): bool\\
			+ exportar\_clave\_publica(): bytes
		};

		\node[cls] (SV) at (5.5, 4.0) {
			\textbf{ServidorVotacionPQC}\\
			\rule{5.2cm}{0.3pt}\\
			- ek\_ac: bytes [1184]\\
			- sk\_sv: bytes [4032]\\
			- pk\_sv: bytes [1952]\\
			- pk\_ar: bytes [1952]\\
			- anfora\_digital: List[BoletaSufragio]\\
			\rule{5.2cm}{0.3pt}\\
			+ recepcionar\_voto(paquete: PaqueteVoto): ReciboVotacion\\
			+ verificar\_credencial(cred: CredencialVotante): bool\\
			+ verificar\_firma\_boleta(blt: BoletaSufragio, sig: bytes): bool\\
			+ clausurar\_urna(): List[BoletaSufragio]\\
			+ emitir\_recibo(blt: BoletaSufragio): ReciboVotacion
		};

		\node[cls] (AC) at (5.5, -0.6) {
			\textbf{AutoridadConteoPQC}\\
			\rule{5.2cm}{0.3pt}\\
			- dk\_ac: bytes [2400]\\
			- ek\_ac: bytes [1184]\\
			- pk\_ar: bytes [1952]\\
			\rule{5.2cm}{0.3pt}\\
			+ generar\_par\_kem(): Tuple[bytes, bytes]\\
			+ exportar\_ek(): bytes\\
			+ escrutinio(anfora: List[BoletaSufragio]): ResultadoEscrutinio
		};

		\node[cls] (Pkt) at (-0.5, 4.0) {
			«contentClass»\\
			\textbf{PaqueteVoto}\\
			\rule{4.4cm}{0.3pt}\\
			+ boleta: BoletaSufragio\\
			+ credencial: CredencialVotante\\
			+ firma\_votante: bytes [3309]\\
			\rule{4.4cm}{0.3pt}\\
			+ serializar\_canonico(): bytes
		};

		\node[cls] (Cred) at (-0.5, 1.3) {
			«contentClass»\\
			\textbf{CredencialVotante}\\
			\rule{4.4cm}{0.3pt}\\
			+ id\_votante: str\\
			+ pk\_votante: bytes [1952]\\
			+ timestamp\_emision: int\\
			+ firma\_ar: bytes [3309]\\
			\rule{4.4cm}{0.3pt}\\
			+ obtener\_bytes(): bytes
		};

		\node[cls] (Blt) at (-0.5, -1.7) {
			«contentClass»\\
			\textbf{BoletaSufragio}\\
			\rule{4.4cm}{0.3pt}\\
			+ id\_eleccion: str\\
			+ ciphertext\_kem: bytes [1088]\\
			+ voto\_cifrado: bytes [48]\\
			+ nonce\_aead: bytes [12]\\
			\rule{4.4cm}{0.3pt}\\
			+ obtener\_bytes(): bytes
		};

		\node[cls] (Rec) at (5.5, -4.6) {
			«contentClass»\\
			\textbf{ReciboVotacion}\\
			\rule{4.6cm}{0.3pt}\\
			+ hash\_boleta: bytes [32]\\
			+ timestamp\_recibo: int\\
			+ firma\_servidor: bytes [3309]\\
			\rule{4.6cm}{0.3pt}\\
			+ obtener\_bytes(): bytes
		};

		\node[cls] (Res) at (-0.5, -4.6) {
			«contentClass»\\
			\textbf{ResultadoEscrutinio}\\
			\rule{4.4cm}{0.3pt}\\
			+ totales\_opciones: Dict[str, int]\\
			+ total\_votos: int\\
			+ hashes\_auditoria: List[str]\\
			\rule{4.4cm}{0.3pt}\\
			+ exportar\_acta(): dict
		};

		\draw[assoc] (V) -- (Pkt) node[midway, above, font=\tiny\rmfamily] {formula} node[near start, below, font=\tiny\rmfamily] {1} node[near end, below, font=\tiny\rmfamily] {0..1};
		\draw[comp] (Pkt) -- (Cred) node[near start, right, font=\tiny\rmfamily] {1} node[near end, right, font=\tiny\rmfamily] {1};
		\draw[comp] (Pkt.south) to[bend right=20] node[near start, left, font=\tiny\rmfamily] {1} node[near end, left, font=\tiny\rmfamily] {1} (Blt.north);
		\draw[assoc] (AR) -- (V) node[midway, left, font=\tiny\rmfamily] {acredita} node[near start, right, font=\tiny\rmfamily] {1} node[near end, right, font=\tiny\rmfamily] {1..*};
		\draw[assoc] (Pkt) -- (SV) node[midway, above, font=\tiny\rmfamily] {envía a SV} node[near start, below, font=\tiny\rmfamily] {0..*} node[near end, below, font=\tiny\rmfamily] {1};
		\draw[assoc] (SV.west) to[bend left=15] node[near start, above, font=\tiny\rmfamily] {1} node[near end, above, font=\tiny\rmfamily] {0..* custodiada} (Blt.east);
		\draw[assoc] (SV) -- (Rec) node[midway, right, font=\tiny\rmfamily] {emite} node[near start, left, font=\tiny\rmfamily] {1} node[near end, left, font=\tiny\rmfamily] {0..*};
		\draw[assoc] (SV) -- (AC) node[midway, right, font=\tiny\rmfamily] {lote mezclado} node[near start, left, font=\tiny\rmfamily] {1} node[near end, left, font=\tiny\rmfamily] {1};
		\draw[assoc] (AC) -- (Res) node[midway, below, font=\tiny\rmfamily] {genera acta} node[near start, above, font=\tiny\rmfamily] {1} node[near end, above, font=\tiny\rmfamily] {1};
	\end{tikzpicture}
	}%
	\caption{Diagrama de clases del Sistema.}
	\label{fig:uml_clases_detallado}
	\par\vspace{3pt}{\footnotesize\textit{Fuente: Elaboración propia .}}
\end{figure}

Las cinco clases de contenido («contentClass») cumplen funciones específicas y verificables: \textbf{CredencialVotante} almacena la clave pública ML-DSA-65 del elector (1952 B) avalada por la firma de la Autoridad de Registro (3309 B); \textbf{BoletaSufragio} encapsula de forma anónima el criptograma ML-KEM-768 (1088 B), el voto cifrado (48 B) y el vector de inicialización (12 B); \textbf{PaqueteVoto} agrupa mediante composición la boleta y la credencial junto con la firma del votante (3309 B) para su transmisión; \textbf{ReciboVotacion} contiene el resumen SHA3-256 (32 B) firmado por el servidor; y \textbf{ResultadoEscrutinio} consolida el cómputo final y la lista de resúmenes auditables. Para que la firma digital produzca siempre el mismo resultado binario, todos los objetos se transforman previamente mediante una función de serialización canónica que ordena alfabéticamente sus atributos sin espacios intermedios y aplica separación de contexto según el tipo de documento (credencial, boleta o recibo).

\subsection{Modelo Entidad-Relación y Privacidad del Sufragio}

Si la tabla de boletas almacenara una clave foránea hacia la cédula del ciudadano, cualquier administrador de base de datos podría cruzar ambas tablas y vulnerar el secreto del sufragio. Para evitarlo, el diseño aplica la \textbf{desvinculación criptográfica de identidad}, ilustrado en el Modelo Entidad-Relación de la Figura \ref{fig:bbdd_er}: al validar el paquete de voto, el Servidor de Votación marca el indicador de sufragio emitido en la rama del padrón (izquierda) y deposita únicamente la boleta cifrada anónima en la rama de la urna digital (derecha), sin relación foránea entre el elector y la boleta.

\begin{figure}[!htbp]
	\centering
	\resizebox{\textwidth}{!}{%
	\begin{tikzpicture}[
		entity/.style={rectangle, draw=PQCMainBlue, fill=BoxBgBlue, line width=1.0pt, font=\scriptsize\bfseries, minimum width=2.9cm, minimum height=0.75cm, align=center},
		rel/.style={diamond, draw=PQCTeal, fill=BoxBgTeal, line width=0.8pt, aspect=1.8, font=\tiny\bfseries, align=center, inner sep=2pt},
		link/.style={line width=0.6pt, draw=black!75},
		attr/.style={ellipse, draw=black!60, fill=white, line width=0.5pt, font=\tiny, align=center, inner sep=1.8pt},
		keyattr/.style={ellipse, draw=PQCMainBlue, fill=BoxBgBlue, line width=1.2pt, font=\tiny\bfseries, align=center, inner sep=2.4pt}
	]
		% --- Cuadro de Simbología / Leyenda ---
		\draw[draw=gray!50, fill=gray!4, rounded corners=3pt, line width=0.5pt] (-10.4, 3.6) rectangle (-6.2, 5.3);
		\node[font=\scriptsize\bfseries\color{PQCDark}] at (-8.3, 5.05) {Simbología};
		\node[keyattr, scale=0.85] (leg_pk) at (-9.4, 4.45) {\underline{id\_pk}};
		\node[right=2pt of leg_pk, font=\tiny\bfseries] {Clave Primaria (PK)};
		\node[attr, scale=0.85] (leg_at) at (-9.4, 3.95) {atributo};
		\node[right=2pt of leg_at, font=\tiny] {Atributo común};

		% --- Entidades Principales ---
		\node[entity] (elec) at (0, 3.8) {ELECCIÓN};
		\node[entity] (padron) at (-6.2, 0.8) {PADRÓN\_VOTANTE};
		\node[entity] (cred) at (-6.2, -3.2) {CREDENCIAL\_AR};
		\node[entity] (anfora) at (6.2, 0.8) {BOLETA\_CIFRADA};
		\node[entity] (recibo) at (6.2, -3.2) {RECIBO\_VOTACIÓN};
		\node[entity] (audit) at (0, -4.6) {BITÁCORA\_AUDITORÍA};
		
		% --- Relaciones ---
		\node[rel] (r1) at (-3.1, 2.3) {Habilita};
		\node[rel] (r2) at (-6.2, -1.2) {Posee};
		\node[rel] (r3) at (3.1, 2.3) {Custodia};
		\node[rel] (r4) at (6.2, -1.2) {Genera};
		\node[rel] (r5) at (0, -0.4) {Registra};
		
		% --- Conexiones y Cardinalidades ---
		\draw[link] (elec) -- (r1) node[midway, above, font=\tiny] {1};
		\draw[link] (r1) -- (padron) node[midway, above, font=\tiny] {N};
		\draw[link] (padron) -- (r2) node[midway, left, font=\tiny] {1};
		\draw[link] (r2) -- (cred) node[midway, left, font=\tiny] {1};
		\draw[link] (elec) -- (r3) node[midway, above, font=\tiny] {1};
		\draw[link] (r3) -- (anfora) node[midway, above, font=\tiny] {N};
		\draw[link] (anfora) -- (r4) node[midway, right, font=\tiny] {1};
		\draw[link] (r4) -- (recibo) node[midway, right, font=\tiny] {1};
		\draw[link] (elec) -- (r5) node[midway, left, font=\tiny] {1};
		\draw[link] (r5) -- (audit) node[midway, left, font=\tiny] {N};
		
		% --- Atributos de ELECCIÓN ---
		\node[keyattr] (k_el) at (0, 5.3) {\underline{id\_eleccion}}; \draw[link] (elec) -- (k_el);
		\node[attr] (a_el_nom) at (-2.3, 5.0) {nombre}; \draw[link] (elec) -- (a_el_nom);
		\node[attr] (a_el_fas) at (2.3, 5.0) {estado\_fase}; \draw[link] (elec) -- (a_el_fas);
		\node[attr] (a_el_ini) at (-4.3, 4.4) {fecha\_inicio}; \draw[link] (elec) -- (a_el_ini);
		\node[attr] (a_el_fin) at (4.3, 4.4) {fecha\_fin}; \draw[link] (elec) -- (a_el_fin);
		\node[attr] (a_el_pkar) at (-3.9, 3.4) {pk\_ar}; \draw[link] (elec) -- (a_el_pkar);
		\node[attr] (a_el_ekac) at (3.9, 3.4) {ek\_ac}; \draw[link] (elec) -- (a_el_ekac);
		
		% --- Atributos de PADRÓN_VOTANTE ---
		\node[keyattr] (k_pa) at (-6.2, 2.3) {\underline{id\_eleccion, id\_votante}}; \draw[link] (padron) -- (k_pa);
		\node[attr] (a_pa_hav) at (-4.0, 1.6) {ha\_votado}; \draw[link] (padron) -- (a_pa_hav);
		\node[attr] (a_pa_tsv) at (-4.0, 0.0) {timestamp\_voto}; \draw[link] (padron) -- (a_pa_tsv);
		\node[attr] (a_pa_nom) at (-8.8, 1.8) {nombres}; \draw[link] (padron) -- (a_pa_nom);
		\node[attr] (a_pa_ape) at (-9.2, 0.8) {apellidos}; \draw[link] (padron) -- (a_pa_ape);
		\node[attr] (a_pa_car) at (-8.8, -0.2) {carrera}; \draw[link] (padron) -- (a_pa_car);
		
		% --- Atributos de CREDENCIAL_AR ---
		\node[keyattr] (k_cr) at (-8.8, -2.4) {\underline{id\_credencial}}; \draw[link] (cred) -- (k_cr);
		\node[attr] (a_cr_idv) at (-9.3, -3.3) {id\_votante (FK)}; \draw[link] (cred) -- (a_cr_idv);
		\node[attr] (a_cr_pkv) at (-8.4, -4.4) {pk\_votante\_hex}; \draw[link] (cred) -- (a_cr_pkv);
		\node[attr] (a_cr_far) at (-6.2, -4.6) {firma\_ar\_hex}; \draw[link] (cred) -- (a_cr_far);
		\node[attr] (a_cr_tse) at (-4.3, -4.2) {timestamp\_emision}; \draw[link] (cred) -- (a_cr_tse);
		
		% --- Atributos de BOLETA_CIFRADA ---
		\node[keyattr] (k_an) at (6.2, 2.3) {\underline{id\_boleta}}; \draw[link] (anfora) -- (k_an);
		\node[attr] (a_an_ide) at (4.0, 1.6) {id\_eleccion (FK)}; \draw[link] (anfora) -- (a_an_ide);
		\node[attr] (a_an_hsh) at (4.0, 0.0) {hash\_sha3}; \draw[link] (anfora) -- (a_an_hsh);
		\node[attr] (a_an_kem) at (9.4, 1.8) {ciphertext\_kem}; \draw[link] (anfora) -- (a_an_kem);
		\node[attr] (a_an_vot) at (9.6, 0.9) {voto\_cifrado}; \draw[link] (anfora) -- (a_an_vot);
		\node[attr] (a_an_non) at (9.4, 0.0) {nonce\_aead}; \draw[link] (anfora) -- (a_an_non);
		\node[attr] (a_an_tsr) at (8.6, -0.9) {timestamp\_recepcion}; \draw[link] (anfora) -- (a_an_tsr);
		
		% --- Atributos de RECIBO_VOTACIÓN ---
		\node[keyattr] (k_re) at (8.8, -2.4) {\underline{id\_recibo}}; \draw[link] (recibo) -- (k_re);
		\node[attr] (a_re_idb) at (9.3, -3.3) {id\_boleta (FK)}; \draw[link] (recibo) -- (a_re_idb);
		\node[attr] (a_re_tsr) at (8.4, -4.4) {timestamp\_recibo}; \draw[link] (recibo) -- (a_re_tsr);
		\node[attr] (a_re_fsv) at (6.2, -4.6) {firma\_sv\_hex}; \draw[link] (recibo) -- (a_re_fsv);
		
		% --- Atributos de BITÁCORA_AUDITORÍA ---
		\node[keyattr] (k_au) at (0, -6.0) {\underline{id\_evento}}; \draw[link] (audit) -- (k_au);
		\node[attr] (a_au_ide) at (-2.7, -5.6) {id\_eleccion (FK)}; \draw[link] (audit) -- (a_au_ide);
		\node[attr] (a_au_tip) at (-2.8, -3.7) {tipo\_evento}; \draw[link] (audit) -- (a_au_tip);
		\node[attr] (a_au_des) at (2.8, -3.7) {descripcion}; \draw[link] (audit) -- (a_au_des);
		\node[attr] (a_au_tse) at (2.7, -5.6) {timestamp\_evento}; \draw[link] (audit) -- (a_au_tse);
	\end{tikzpicture}
	}%
	\caption{Modelo Entidad-Relación de la base de datos.}
	\label{fig:bbdd_er}
	\par\vspace{3pt}{\footnotesize\textit{Fuente: Elaboración propia.}}
\end{figure}

En la Tabla \ref{tab:diccionario_datos} se sintetiza la estructura de las seis tablas relacionales normalizadas en Tercera Forma Normal (3FN), detallando sus llaves primarias, atributos principales y mecanismos de seguridad.

\begin{table}[!htbp]
	\centering
	\renewcommand{\arraystretch}{1.22}
	\footnotesize
	\caption{Esquema relacional en tercera forma normal.}
	\label{tab:diccionario_datos}
	\begin{tabularx}{\textwidth}{|>{\bfseries\raggedright\arraybackslash}p{3.0cm}|>{\raggedright\arraybackslash}p{2.7cm}|>{\raggedright\arraybackslash}p{4.3cm}|>{\raggedright\arraybackslash}X|}
		\hline
		\textbf{Tabla Relacional} & \textbf{Clave Primaria} & \textbf{Atributos Principales} & \textbf{Función y Regla de Seguridad} \\
		\hline
		1. Elecciones & id\_eleccion & nombre, fecha\_inicio, fecha\_fin, estado\_fase, pk\_ar, ek\_ac & Almacena los parámetros del comicio y las claves públicas de las autoridades. \\
		\hline
		2. Padrón de Votantes & (id\_eleccion, id\_votante) & nombres, apellidos, carrera, ha\_votado, timestamp\_voto & Controla la elegibilidad y el sufragio único en tiempo $\mathcal{O}(\log n)$ sin vínculo hacia la boleta. \\
		\hline
		3. Credenciales AR & id\_credencial & id\_votante (FK única), pk\_votante\_hex, firma\_ar\_hex & Registra el certificado digital ML-DSA-65 (3309 B) emitido al empadronar al ciudadano. \\
		\hline
		4. Ánfora de Boletas & id\_boleta & id\_eleccion (FK), ciphertext\_kem, voto\_cifrado, nonce, hash\_sha3 & Custodia las boletas anónimas con índice único sobre el resumen SHA3-256 (64 caracteres). \\
		\hline
		5. Recibos de Votación & id\_recibo & id\_boleta (FK única), timestamp\_recibo, firma\_sv\_hex & Conserva el comprobante firmado por el Servidor de Votación para auditoría individual. \\
		\hline
		6. Bitácora de Eventos & id\_evento & id\_eleccion (FK), tipo\_evento, descripcion, timestamp\_evento & Registra los eventos de apertura, cierre y escrutinio excluyendo direcciones IP de votantes. \\
		\hline
	\end{tabularx}
	\par\vspace{3pt}{\footnotesize\textit{Fuente: Elaboración propia.}}
\end{table}
